\documentclass[11pt,a4paper]{article}
\usepackage[margin=2.32cm]{geometry}
\usepackage[T1]{fontenc}
\usepackage{lmodern,microtype}
\usepackage{xcolor}
\usepackage{amsmath,amssymb,amsthm,mathtools}
\usepackage{hyperref}
\hypersetup{colorlinks=true,linkcolor=black,urlcolor=blue!50!black,citecolor=black}
\usepackage{enumitem}
\setlength{\parskip}{1.5pt}
\setlength{\emergencystretch}{2em}
\newtheorem{theorem}{Theorem}[section]
\newtheorem{lemma}[theorem]{Lemma}
\newtheorem{proposition}[theorem]{Proposition}
\newtheorem{corollary}[theorem]{Corollary}
\theoremstyle{remark}\newtheorem{remark}[theorem]{Remark}
\newcommand{\reg}{\operatorname{reg}}
\newcommand{\pd}{\operatorname{pd}}
\newcommand{\GG}{\mathcal G}
\newcommand{\Bad}{\mathcal B}
\newcommand{\Good}{\mathcal C}
\newcommand{\pol}{\mathrm{pol}}
\title{\textbf{Linear resolutions of powers of edge ideals of gap-free graphs with matching number four}}
\author{Luis Alfredo Dupont Garc\'ia\\[2pt]
\small Facultad de Matem\'aticas, Universidad Veracruzana\\
\small Paseo No.~112, Lote 2, Col.~Nuevo Xalapa, C.P.~91097, Xalapa, Veracruz, Mexico\\
\small \texttt{ldupont@uv.mx}\quad ORCID: \href{https://orcid.org/0000-0002-9044-3547}{0000-0002-9044-3547}}
\date{}
\begin{document}
\maketitle
\begin{abstract}
For a finite simple gap-free graph $G$, let $\Bad(G)$ denote the vertices $v$ for which $G-N_G[v]$ is not co-chordal, and put $\Good(G)=V(G)\setminus\Bad(G)$. We prove a sufficient condition for all ordinary powers of its edge ideal to have linear resolutions from the second power onward: every induced anticycle of length at least five must contain an edge with both endpoints in $\Good(G)$. The argument uses variable generators of augmented colon ideals to eliminate anticycles that may persist in the ordinary quadratic colons. It yields linear resolutions for all powers $I(G)^q$, $q\ge2$, whenever the ordinary matching number is at most four, which answers the first part of a question of Erey, H\`a, Hibi, Kara and Morey. The criterion also applies to certain graphs with adjacent exceptional vertices and matching number five. The proofs are theoretical and valid over arbitrary fields. We do not assert linear quotients.
\end{abstract}
\noindent\textbf{2020 Mathematics Subject Classification.} Primary 13D02; Secondary 05E40, 13F55.\par
\noindent\textbf{Keywords.} Edge ideals; Castelnuovo--Mumford regularity; linear powers; gap-free graphs; matching number; induced anticycles; augmented colon ideals.\par


\section{Introduction and statements}
Let $G$ be a finite simple graph with $E(G)\ne\varnothing$, let $K$ be any field, and set
\[
 S=K[x_v:v\in V(G)],\qquad I=I(G)=(x_ux_v:uv\in E(G)).
\]
We write $\reg I$ for the Castelnuovo--Mumford regularity of the ideal as a graded $S$-module. Since $I^q$ is equigenerated in degree $2q$, it has a linear resolution if and only if $\reg I^q=2q$.
A graph is \emph{gap-free} if it contains no induced $2K_2$, and is \emph{co-chordal} if its complement is chordal. An \emph{anticycle} of length $r$ means an induced copy of $\overline{C_r}$. Write $N_G[v]$ for the closed neighbourhood of $v$ and define
\[
 \Bad(G)=\{v\in V(G):G-N_G[v]\text{ is not co-chordal}\},\qquad
 \Good(G)=V(G)\setminus\Bad(G).
\]
We call the vertices of $\Bad(G)$ \emph{exceptional}. Write $\mu(G)$ for the ordinary matching number, as opposed to the induced matching number. The following is the central result of this work.

\begin{theorem}[Anticycle criterion]\label{thm:main}
Let $K$ be any field. Suppose that $G$ is gap-free and that
\begin{equation}\tag{A}\label{eq:A}
  E\bigl(G[W\cap\Good(G)]\bigr)\ne\varnothing
  \qquad\text{for every induced anticycle $G[W]\cong\overline{C_r}$, $r\ge5$.}
\end{equation}
Then $\reg I(G)\le3$ and
\[
 \boxed{\reg I(G)^q=2q\qquad(q\ge2)}.
\]
\end{theorem}

\begin{corollary}\label{cor:matching}
If $G$ is gap-free and $\mu(G)\le4$, then $\reg I(G)^q=2q$ for every $q\ge2$ and every field $K$.
\end{corollary}

Nevo and Peeva \cite{NP} conjectured that gap-free graphs with $\reg I(G)=3$ have linear powers from the square onward; the more general eventual-linearity conjecture for gap-free graphs is still open. Previous sufficient conditions include gap-free and cricket-free graphs \cite{Banerjee}, gap-free and diamond-free graphs \cite{EreyDiamond}, and graphs with neither an induced $C_4$ nor an induced $2K_2$ \cite{EreyC4}. Banerjee--Beyarslan--H\`a \cite[Theorem 4.5]{BBH} proved the conclusion for gap-free graphs that are locally linear, that is, under the hypothesis $\Bad(G)=\varnothing$. In that case \eqref{eq:A} holds trivially, since every anticycle has edges and all its vertices lie in $\Good(G)=V(G)$; thus Theorem~\ref{thm:main} contains \cite[Theorem 4.5]{BBH}. Recent work also investigates the stronger property of linear quotients, including graph reductions and preservation operations \cite{EFFHKM,EHHKM}.

Minh--Vu \cite{MV} proved that $I(G)^2$ has a linear resolution if and only if $G$ is gap-free and $\reg I(G)\le3$, and that $I(G)^3$ has a linear resolution if and only if $G$ is gap-free and $\reg I(G)\le4$. Hence, once $G$ is gap-free and $\reg I(G)\le3$, the cases $q=2,3$ are already covered, and the uniform content here concerns $q\ge4$ in the presence of exceptional vertices. In a preprint dated October 15, 2026, Erey--H\`a--Hibi--Kara--Morey proved that all powers $I(G)^q$, $q\ge2$, of a gap-free graph with matching number three have linear quotients, and \cite[Question 5.1]{EHHKM} explicitly asked whether all powers beginning at the square of a gap-free graph with matching number four have linear resolutions, and also whether they have linear quotients. Corollary~\ref{cor:matching} answers the first question affirmatively; the second remains open. The broader anticycle condition~\eqref{eq:A} also applies when exceptional vertices are adjacent.

The mechanism is to \emph{retain}, rather than discard, the variable part of the augmented colons $(I^{s+1},m_1,\dots,m_{j-1}):m_j$. These variables eliminate anticycles that can survive in the quadratic part $(I^{s+1}:m_j)$. We give a direct alternating-trail proof that permits an ordinary monomial ordering of the $m_j$, so that the decisive comparison becomes an exchange of an exceptional variable for a non-exceptional one.

Section~\ref{sec:colons} develops the alternating-trail exchange and the structure of augmented colons; Section~\ref{sec:anticycles} studies anticycles surviving in edge colons; Section~\ref{sec:proof} proves Theorem~\ref{thm:main} and gives a polynomial-time test for \eqref{eq:A}; Section~\ref{sec:example} proves Corollary~\ref{cor:matching} and analyzes a ten-vertex example with adjacent exceptional vertices. We use Banerjee's published persistence lemmas for anticycles in polarized edge colons \cite[Lemmas 6.14 and 6.15]{Banerjee}. No computational calculation is a premise of the results.

\section{Alternating trails and augmented colons}\label{sec:colons}
For monomials $u,v$ we abbreviate $u:v=u/\gcd(u,v)$; for ideals, $(J:v)$ is the standard colon ideal. Write $\GG(J)$ for the minimal monomial generating set of a monomial ideal $J$.

\begin{lemma}[Alternating-trail exchange]\label{lem:exchange}
Let $u,v\in\GG(I^s)$ with $u\ne v$, and put $Q_v=(I^{s+1}:v)$. If $u:v\notin Q_v$, there exist vertices $a_1,b_1,\dots,a_h,b_h$, with $h\ge1$, such that
\[
 u:v=\prod_{i=1}^h x_{a_i},\qquad v:u=\prod_{i=1}^h x_{b_i},
 \qquad w_i=\frac{x_{a_i}v}{x_{b_i}}\in\GG(I^s),\quad w_i:v=x_{a_i}.
\]
In particular, for any monomial ordering $>$, if $u>v$, some $w_i>v$ and $w_i:v=x_{a_i}\mid u:v$.
\end{lemma}
\begin{proof}
Choose factorizations of $u$ and $v$ into $s$ edge occurrences, coloured red and blue, respectively, and cancel identical occurrences of opposite colours. At each vertex, pair as many red incidences with blue incidences as possible. Following edges and pairings decomposes the remaining edge occurrences into a finite collection of alternating trails. The unpaired incidences are the end-incidences of the open trails; at each vertex they all have the same colour, and their number is the positive (red) or negative (blue) exponent difference between $u$ and $v$ at that vertex.

An alternating trail with two red endpoints $a,b$ has one more red than blue edge occurrence. Writing $R_T$ and $B_T$ for the products of its red and blue edges gives $R_T=x_ax_bB_T$. Replacing its blue occurrences in the factorization of $v$ by its red occurrences shows
\[
 x_ax_bv=(v/B_T)R_T\in I^{s+1}.
\]
Since $x_ax_b\mid u:v$, this would imply $u:v\in Q_v$, a contradiction. Since the total numbers of red and blue edge occurrences are equal, there can be no blue--blue open trail either. All open trails therefore have one red and one blue endpoint, and have equal numbers of red and blue occurrences. Swapping those occurrences yields $w_i=x_{a_i}v/x_{b_i}\in I^s$, necessarily minimal because its degree is $2s$. A vertex cannot carry unpaired incidences of both colours, so $a_i\ne b_i$, and therefore $w_i:v=x_{a_i}$. Accounting for all unpaired incidences gives the displayed factorizations. Since $u\ne v$, at least one open trail exists.

For the last assertion, if every $x_{a_i}<x_{b_i}$, multiplicativity of the monomial order would imply $u:v<v:u$, hence $u<v$, a contradiction. Choose an $i$ with $x_{a_i}>x_{b_i}$; then $w_i>v$.
\end{proof}

\begin{proposition}[Quadratic and augmented colons]\label{prop:colons}
For every graph $G$, every $s\ge1$, and every $m\in\GG(I^s)$, the ideal $Q_m=(I^{s+1}:m)$ is generated by quadratic monomials. Furthermore, let $m_1>\cdots>m_t$ be the minimal generators of $I^s$ in decreasing order for \emph{any} monomial ordering, put
\[
 D_j=(I^{s+1},m_1,\ldots,m_{j-1}):m_j
 \qquad\text{and}\qquad
 Z_j=\{v\in V(G):x_v\in D_j\}.
\]
Then
\begin{equation}\label{eq:D}
 D_j=Q_{m_j}+(x_v:v\in Z_j).
\end{equation}
\end{proposition}
\begin{proof}
The quadratic generation is Banerjee's Theorem~6.1 \cite{Banerjee}; it also follows directly from the excess-red-edge variant of the preceding proof: if $f m\in I^{s+1}$, compare a product of $s+1$ red edges dividing $fm$ with the $s$ blue edges of $m$. A red--red alternating trail must exist by the excess of one red edge. Its endpoints $a,b$ satisfy $x_ax_bm\in I^{s+1}$ and $x_ax_b\mid f$. There are no generators of degree less than two.

For the second assertion, the inclusion $\supseteq$ is clear. Conversely, the monomial ideal $D_j$ is generated by the generators of $Q_{m_j}$ together with the monomials $m_i:m_j$ for $i<j$. Each of the latter either lies in $Q_{m_j}$ or, by Lemma~\ref{lem:exchange}, is divisible by a variable $x_a=w:m_j$, where $w>m_j$ is another minimal generator of $I^s$, say $w=m_k$ with $k<j$. In the second case $x_a\in D_j$, that is, $a\in Z_j$. This proves \eqref{eq:D}.
\end{proof}

Banerjee's Theorem~4.12 \cite{Banerjee} provides such an exchange for a special ordering based on maximal edge expressions. Proposition~\ref{prop:colons} allows a monomial ordering of the original variables. It does \emph{not} assert that $I^s$ has linear quotients: the quadratic part $Q_{m_j}$ remains.

\section{Anticycles in edge colons}\label{sec:anticycles}
Let $m$ be a product of $s$ edges and let $H_m$ be the graph of the polarization of $Q_m=(I^{s+1}:m)$. We use the following published facts:
\begin{equation}\label{eq:Ban}
\begin{gathered}
 G\text{ gap-free}\ \Longrightarrow\ H_m\text{ gap-free},\\
 G\text{ gap-free and } H_m[W]\cong\overline{C_r},\ r\ge5\ \Longrightarrow\
 W\subseteq V(G),\quad G[W]=H_m[W].
\end{gathered}
\end{equation}
The first statement is Banerjee's Lemma~6.14 and the second is his Lemma~6.15 \cite{Banerjee}. New polarization vertices are leaves \cite[Lemma 6.11]{Banerjee}, hence never belong to the relevant anticycles.

\begin{lemma}[An edge facing a surviving anticycle]\label{lem:anticycle}
Suppose that $G$ is gap-free and $H_m[W]\cong\overline{C_r}$ for some $r\ge5$. For every edge $ab$ in any factorization of $m$, both endpoints lie outside $W$. Exactly one endpoint is anticomplete to $W$, while the neighbourhood in $W$ of the other endpoint is a vertex cover of $G[W]$.
\end{lemma}
\begin{proof}
By \eqref{eq:Ban} the two graphs induce the same anticycle on $W$. Label $W=\{w_0,\ldots,w_{r-1}\}$ cyclically so that $w_iw_{i+1}$ is a nonedge, with indices modulo $r$. If $x\in N_G(a)$ and $y\in N_G(b)$, then
\begin{equation}\label{eq:cross}
 xy\,m=(xa)(by)\frac{m}{ab}\in I^{s+1},
\end{equation}
so $xy$ is an edge of $H_m$ when $x\ne y$.

Suppose $a=w_0$. The edge $w_{r-1}w_1$ belongs to $G$ and both its endpoints are non-neighbours of $a$. Since $b$ is a neighbour of $a=w_0$ whereas $w_1,w_{r-1}$ are not, $b\notin\{w_1,w_{r-1}\}$. As $G$ is gap-free, $b$ is adjacent to $w_1$ or $w_{r-1}$. Using $w_2\in N_G(a)$ or $w_{r-2}\in N_G(a)$ in \eqref{eq:cross}, respectively, adds the forbidden edge $w_1w_2$ or $w_{r-2}w_{r-1}$ to $H_m$. Therefore $a\notin W$, and similarly $b\notin W$. (This step parallels an argument in the proof of \cite[Theorem 6.17]{Banerjee}.)

Put $X=N_G(a)\cap W$, $Y=N_G(b)\cap W$ and $R=W\setminus(X\cup Y)$. Equation~\eqref{eq:cross} forbids consecutive positions of the complementary cycle with one in $X$ and the other in $Y$. Every $w_i\in X$ has at least one of its two consecutive positions in $X$: otherwise $aw_i$ and $w_{i-1}w_{i+1}$ would form an induced $2K_2$; the same statement holds for $Y$. Thus $X\cap Y=\varnothing$, since a vertex of the intersection would have a consecutive neighbour in $X$, forbidden by its membership in $Y$.

Any edge in $G[R]$ would form a gap with $ab$. Thus $R$ is independent in $G[W]=\overline{C_r}$, hence has at most two vertices and, if two, they are consecutive on $C_r$. The cycle $C_r-R$ is therefore connected. Since it contains no edge from $X$ to $Y$, one of $X,Y$ must be empty. They cannot both be empty because $G[W]$ has edges. The endpoint with empty neighbourhood is anticomplete to $W$. The neighbourhood of the other endpoint is $W\setminus R$, which meets every edge of $G[W]$ because $R$ is independent.
\end{proof}

\section{Proof of the anticycle criterion}\label{sec:proof}
For an induced subgraph $F$ of $G$ we write $S_F=K[x_v:v\in V(F)]$.

\begin{lemma}[Initial regularity]\label{lem:initial}
If $G[\Bad(G)]$ is co-chordal, then $\reg(S/I(G))\le2$. In particular, if $I(G)\ne0$, then $\reg I(G)\le3$.
\end{lemma}
\begin{proof}
Delete the vertices of $\Good(G)$ one by one. If $F$ is the current induced graph and $c\in\Good(G)\cap V(F)$, then $F-N_F[c]$ is an induced subgraph of the co-chordal graph $G-N_G[c]$, hence co-chordal. Since
\[
 (I(F):x_c)=(x_z:z\in N_F(c))+I(F-N_F[c]),
\]
Fr\"oberg's theorem \cite{Froberg} gives $\reg\big(S_F/(I(F):x_c)\big)\le1$. The standard exact sequence
\[
 0\to S_F/(I(F):x_c)(-1)\xrightarrow{\;\cdot x_c\;}S_F/I(F)\to S_F/(I(F),x_c)\to0
\]
yields
\[
 \reg(S_F/I(F))\le\max\{2,\reg(S_{F-c}/I(F-c))\}.
\]
After all deletions, only $G[\Bad(G)]$ remains, whose edge-ideal quotient has regularity at most one. Hence the claimed bound follows.
\end{proof}

\begin{proof}[Proof of Theorem~\ref{thm:main}]
Write $B=\Bad(G)$ and $C=\Good(G)$. Since $G$ is gap-free, $G[B]$ has no induced anticycle of length four. If it contained one of length $r\ge5$, its vertices $W\subseteq B$ would violate \eqref{eq:A}. Thus $G[B]$ is co-chordal and Lemma~\ref{lem:initial} gives $\reg I\le3$.

Choose a monomial order with $x_c>x_b$ for every $c\in C$ and $b\in B$; a lexicographic order listing all variables of $C$ first will do. For each $s\ge1$, order $\GG(I^s)$ decreasingly as $m_1>\cdots>m_t$ and use the ideals $Q_{m_j}$, $D_j$ and the sets $Z_j$ of Proposition~\ref{prop:colons}. We claim
\begin{equation}\label{eq:boundD}\reg D_j\le2\quad\text{for every }j.\end{equation}
Let $Q'_j$ be the ideal generated by those minimal generators of $Q_{m_j}$ involving no variable $x_z$ with $z\in Z_j$. Then
\[
 D_j=Q'_j+(x_z:z\in Z_j).
\]
The variables $x_z$, $z\in Z_j$, form a regular sequence on $S/Q'_j$, since $Q'_j$ involves only the remaining variables. Thus $\reg(S/D_j)=\reg(S'/Q'_j)$ with $S'=K[x_v:v\notin Z_j]$, and
\[
 \reg D_j=\max\{1,\reg Q'_j\},
\]
with $\reg(0)=-\infty$. Up to isolated polarization vertices, the polarization of $Q'_j$ is the edge ideal of $H_{m_j}-Z_j$. In particular, if $\reg D_j>2$, then by Fr\"oberg's theorem that residual polarized graph is not co-chordal and, being an induced subgraph of the gap-free graph $H_{m_j}$, contains an induced anticycle $W$ of length $r\ge5$. Polarization leaves cannot belong to $W$, and $W\cap Z_j=\varnothing$. Lemma~\ref{lem:anticycle} applied to an edge in a factorization of $m_j$ gives an edge $bp$ with $b\notin W$ anticomplete to $W$, and with $N_G(p)\cap W$ a vertex cover of $G[W]$. Since $G-N_G[b]$ contains this anticycle, $b\in B$. We make \emph{no} assumption that $p\in C$ or that $B$ is independent.

By \eqref{eq:A}, some edge $e$ of $G$ has both endpoints in $W\cap C$. The vertex cover $N_G(p)\cap W$ meets $e$, at a vertex $w\in C$. Thus $pw\in E(G)$ and
\[
 m'=\frac{x_wm_j}{x_b}=\frac{m_j}{x_bx_p}(x_px_w)\in\GG(I^s).
\]
As $x_w>x_b$, we have $m'>m_j$, so $m'=m_k$ for some $k<j$; moreover $m':m_j=x_w$ because $w\ne b$. Hence $x_w\in D_j$, that is, $w\in Z_j$. This contradicts $w\in W$ and $W\cap Z_j=\varnothing$. Consequently \eqref{eq:boundD} holds.

Finally, set $L_0=I^{s+1}$ and $L_j=L_{j-1}+(m_j)$. Since $L_t=I^s$, the successive quotients $L_j/L_{j-1}\cong S/D_j(-2s)$ form a filtration of $I^s/I^{s+1}$. By \eqref{eq:boundD} they have regularity at most $2s+1$, hence so does $I^s/I^{s+1}$, and the exact sequence $0\to I^{s+1}\to I^s\to I^s/I^{s+1}\to0$ gives
\[
 \reg I^{s+1}\le\max\{\reg I^s,2s+2\}.
\]
The bound $\reg I\le3$ begins the induction at $s=1$, showing $\reg I^q\le2q$ for every $q\ge2$. Since $I^q$ is generated in degree $2q$, equality holds.
\end{proof}

\begin{proposition}[A polynomial-time local test]\label{prop:localtest}
Suppose that $G$ is gap-free, and put $B=\Bad(G)$ and $C=\Good(G)$. Then condition~\eqref{eq:A} is equivalent to
\begin{equation}\label{eq:localtest}
  G[B\cup T]\text{ is co-chordal for every independent set }
  T\subseteq C\text{ with }|T|\le2.
\end{equation}
In particular, condition~\eqref{eq:A} is decidable by at most
$1+|C|+\binom{|C|}{2}$ co-chordality tests in addition to the $|V(G)|$ local tests determining $B$. Hence it is recognizable in polynomial time.
\end{proposition}
\begin{proof}
If~\eqref{eq:localtest} fails for some independent $T\subseteq C$, then $G[B\cup T]$ is not co-chordal. Its complement contains an induced cycle of length at least four; since $G$ is gap-free, a $4$-cycle is excluded. Thus an induced anticycle $G[W]\cong\overline{C_r}$ with $r\ge5$ lies in $G[B\cup T]$. Its vertices in $C$ lie in the independent set $T$, so $G[W\cap C]$ has no edge, contradicting~\eqref{eq:A}.

Conversely, suppose~\eqref{eq:A} fails for an induced anticycle on $W$. Then $T=W\cap C$ is independent in $G[W]\cong\overline{C_r}$, which has independence number two. Thus $|T|\le2$ and $W\subseteq B\cup T$, so $G[B\cup T]$ is not co-chordal. This contradicts~\eqref{eq:localtest}.

All the relevant sets $T$ have size at most two; both computing $B$ and testing chordality are polynomial-time procedures.
\end{proof}

\begin{remark}\label{rem:complementary}
For $F=\overline G$, the graph $G$ is gap-free exactly when $F$ has no induced $C_4$, and $v\in\Bad(G)$ exactly when $F[N_F(v)]$ is not chordal. When $G$ is gap-free, every hole of $F$ (induced cycle of length at least four) has length at least five, and condition~\eqref{eq:A} says equivalently that every hole of $F$ contains two vertices from $\Good(G)$ that are nonadjacent in $F$.
\end{remark}

\section{Matching number and a ten-vertex example}\label{sec:example}
\begin{lemma}\label{lem:match}
If $G$ is gap-free and two adjacent vertices $u,v$ belong to $\Bad(G)$, then $\mu(G)\ge5$.
\end{lemma}
\begin{proof}
There are induced anticycles $W\subseteq V(G)\setminus N_G[u]$ and $Z\subseteq V(G)\setminus N_G[v]$ of lengths at least five. Since $uv$ is an edge, $u,v\notin W\cup Z$. The intersection $R=W\cap Z$ is independent: an edge inside $R$ would form an induced $2K_2$ with $uv$. Therefore $|R|\le2$.

An anticycle with at least five vertices, after removing at most one prescribed vertex, contains two disjoint edges: choose four consecutive positions on the complementary cycle avoiding that vertex, and take the two short diagonals. If $R=\varnothing$, select two disjoint edges in each of $W,Z$; if $R=\{c\}$, select them in $W\setminus\{c\}$ and $Z$; if $R=\{c,d\}$, select them in $W\setminus\{c\}$ and $Z\setminus\{d\}$. In each case the two supports are disjoint. These four edges, together with $uv$, form a matching of size five.
\end{proof}
\begin{proof}[Proof of Corollary~\ref{cor:matching}]
By Lemma~\ref{lem:match}, $\Bad(G)$ is independent whenever $\mu(G)\le4$. If $W$ is an induced anticycle of length $r\ge5$, then $|W\cap\Bad(G)|\le2$, because the independence number of $\overline{C_r}$ is two. Consequently $W\cap\Good(G)$ has at least three vertices, which cannot be an independent set in $\overline{C_r}$. Condition~\eqref{eq:A} follows; apply Theorem~\ref{thm:main}.
\end{proof}

Define the graph $G_{10}$ on
\[ V=\{u,v,w_0,w_1,w_2,w_3,w_4,z_2,z_3,z_4\} \]
with the following 26 edges:
\begin{align*}
 &w_0w_2,\ w_2w_4,\ w_4w_1,\ w_1w_3,\ w_3w_0,\\
 &w_0z_2,\ z_2w_1,\ w_1z_3,\ z_3z_4,\ z_4w_0,\\
 &uv,\ vw_2,\ vw_3,\ vw_4,\ uz_2,\ uz_3,\ uz_4,\\
 &w_i z_j\quad(i,j\in\{2,3,4\}).
\end{align*}
Let $W=\{w_0,\ldots,w_4\}$ and $Z=\{w_0,w_1,z_2,z_3,z_4\}$.
\begin{proposition}\label{prop:G10}
The graph $G_{10}$ is gap-free, $\Bad(G_{10})=\{u,v\}$, $uv\in E(G_{10})$, and $\mu(G_{10})=5$. It satisfies condition~\eqref{eq:A}, and hence
\[
 \reg I(G_{10})=3,\qquad \reg I(G_{10})^q=2q\quad(q\ge2).
\]
In particular, condition~\eqref{eq:A} allows adjacent exceptional vertices and matching number five.
\end{proposition}
\begin{proof}
The complement $F=\overline{G_{10}}$ is the union of two five-rim wheels, with centres $u$ and $v$, glued precisely along the edge $w_0w_1$. The rims are the cycles
\[
 w_0-w_1-w_2-w_3-w_4-w_0,\qquad
 w_0-w_1-z_4-z_2-z_3-w_0.
\]
There are no additional edges between the non-common parts. An induced cycle of length at least four in a clique-sum along an edge must lie in one summand. Indeed, a cycle meeting both sides passes through both separator vertices; the two paths between them have interior vertices on opposite sides, so the separator vertices are not consecutive on the cycle and their common edge is a chord. In a wheel the induced cycle must avoid the universal centre, and is therefore its rim. Thus the only induced holes of $F$ are $W,Z$, both of length five. In particular $F$ is $C_4$-free, so $G_{10}$ is gap-free.

For every vertex $x$ we have, with the complement taken on the vertex set $N_F(x)$,
\[
 V(G_{10}-N_{G_{10}}[x])=N_F(x),\qquad
 G_{10}-N_{G_{10}}[x]=\overline{F[N_F(x)]}.
\]
Consequently, $x\in\Bad(G_{10})$ if and only if $F[N_F(x)]$ is not chordal. This identifies the exceptional vertices: the neighbourhoods of $u$ and $v$ in $F$ induce the two rim pentagons, while $F[N_F(w_0)]$ and $F[N_F(w_1)]$ are paths on five vertices, and the six other neighbourhood graphs are paths on three vertices. Hence $\Bad(G_{10})=\{u,v\}$. Both holes $W,Z$ avoid that set and so satisfy \eqref{eq:A}. Finally,
\[
 uv,\quad w_1w_3,\quad w_2w_4,\quad w_0z_2,\quad z_3z_4
\]
form a perfect matching, so $\mu(G_{10})=5$. Theorem~\ref{thm:main} gives the power formula. The induced pentagon $G_{10}[W]$ makes $I(G_{10})$ non-linear by Fr\"oberg's theorem, while $\reg I(G_{10})\le3$, proving $\reg I(G_{10})=3$.
\end{proof}

\begin{remark}[Comparison with earlier sufficient criteria]\label{rem:priority}
The graph $G_{10}$ contains an induced cricket, diamond, and $C_4$; it is not locally linear and has $\mu(G_{10})=5$. Thus the standard forbidden-subgraph and local-linearity criteria do not apply \cite{Banerjee,EreyDiamond,EreyC4,BBH}. Among the established base cases for these preservation operations are the co-chordal graphs \cite{HHZ} and the gap-free graphs with ordinary matching number at most three \cite{EHHKM}. A separate finite inverse-operation check, allowing gap-free edge deletions and deletion of open or closed twins (as well as pendant or isolated vertices), yields $16$ isomorphism classes from $G_{10}$; none is co-chordal or has ordinary matching number at most three. The script performing this check is available from the author upon request; the check is \emph{not} part of the proof of Proposition~\ref{prop:G10}, and it does not exclude other sufficient criteria.

An independent Hochster-formula computation over $\mathbb F_2$ gives
\[
 \beta_{1,j}(I(G_{10}))=0\ \ (j\ge4),\qquad \beta_{2,5}(I(G_{10}))=2,\qquad \pd I(G_{10})=6,
\]
where homological degrees are counted in the minimal free resolution of the ideal $I(G_{10})$ itself, not of $S/I(G_{10})$. The two nonlinear Betti numbers $\beta_{2,5}$ come from the two induced pentagons $W$ and $Z$. Thus the resolution of $I(G_{10})$ first fails to be linear in homological degree~$2$, so $\operatorname{index}(I(G_{10}))=2$ in the sense of \cite{BigdeliHZ}, whereas $\pd I(G_{10})=6$. Hence $I(G_{10})$ has neither maximal finite index ($\operatorname{index}=\pd$, \cite{BigdeliHZ}) nor almost maximal finite index ($\operatorname{index}=\pd-1$, \cite{Bigdeli}), and the results on these classes do not cover this example. These auxiliary calculations provide a comparison of scope, not a computer-assisted proof of the linearity of the powers.
\end{remark}

\begin{remark}[Why the order matters]
Write $J=I(G_{10})$. The neighbourhoods of $u,v$ in $G_{10}$ are
\[
 N(u)=\{v,z_2,z_3,z_4\},\qquad N(v)=\{u,w_2,w_3,w_4\}.
\]
All products between the two neighbourhoods already belong to $J$, so $(J^2:uv)=J$. An ordering beginning with $uv$ has first augmented colon $J$, of regularity three. By contrast, if the eight variables of $\Good(G_{10})$ precede $u,v$, the edge $uv$ is last and its augmented colon is
\[
 J+(x_{w_2},x_{w_3},x_{w_4},x_{z_2},x_{z_3},x_{z_4}),
\]
whose residual quadratic part is the edge $uv$. This illustrates the use of the linear components of the augmented colon; it does not assert that every monomial order works.
\end{remark}

\begin{remark}[A boundary example]
Let $H=G_{10}-\{w_2z_4,w_4z_3\}$. The complement of $H$ is obtained from the two-wheel clique-sum by adding the two indicated edges. The graph $H$ remains gap-free, has matching number five and exceptional set $\Bad(H)=\{u,v,w_0,w_1\}$. The set $W'=\{u,v,w_0,w_2,z_4\}$ induces an anticycle in $H$, but the only vertices of $W'\cap\Good(H)$ are $w_2,z_4$, which are not adjacent in $H$. Thus condition~\eqref{eq:A} fails. This does \emph{not} imply that a power of $I(H)$ has a nonlinear resolution: the criterion is sufficient, not necessary.
\end{remark}

\section{Scope and questions}
Theorem~\ref{thm:main} concerns ordinary powers and linear \emph{resolutions}, not linear quotients; in particular, the second part of \cite[Question 5.1]{EHHKM} remains open. The theorem does not settle the general Nevo--Peeva conjecture, nor the matching-number-five case for graphs failing~\eqref{eq:A}, such as the graph $H$ of the boundary example.

The matching-number scale is sharp at the other end. As recalled in \cite[Section 5]{EHHKM}, the complement of the icosahedral graph considered by Nevo and Peeva \cite{NP} is gap-free, has matching number six and edge-ideal regularity four; its square is not linear, while its cube is. Hence no statement asserting linearity of all powers from the square onward can hold for matching number six. Together with Corollary~\ref{cor:matching}, this leaves matching number five as the only intermediate case, and Theorem~\ref{thm:main} covers those graphs of matching number five that satisfy~\eqref{eq:A}, such as $G_{10}$.

The proof of the main theorem uses the published persistence results for anticycles in polarized colons \cite{Banerjee} with their stated hypotheses. Proposition~\ref{prop:localtest} supplies a polynomial-time recognition procedure for condition~\eqref{eq:A}. The finite tests mentioned in Remark~\ref{rem:priority} are optional illustrations and are not used to prove any of the stated regularity formulas.

\begin{thebibliography}{99}\small
\bibitem{Banerjee} A. Banerjee, \emph{The regularity of powers of edge ideals}, J. Algebraic Combin. \textbf{41} (2015), no.~2, 303--321. \url{https://doi.org/10.1007/s10801-014-0537-2}.
\bibitem{BBH} A. Banerjee, S. K. Beyarslan, and H. T. H\`a, \emph{Regularity of powers of edge ideals: from local properties to global bounds}, Algebraic Combin. \textbf{3} (2020), no.~4, 839--854. \url{https://doi.org/10.5802/alco.119}.
\bibitem{Bigdeli} M. Bigdeli, \emph{Edge ideals with almost maximal finite index and their powers}, J. Algebraic Combin. \textbf{54} (2021), 947--978. \url{https://doi.org/10.1007/s10801-021-01032-y}.
\bibitem{BigdeliHZ} M. Bigdeli, J. Herzog, and R. Zaare-Nahandi, \emph{On the index of powers of edge ideals}, Comm. Algebra \textbf{46} (2018), no.~3, 1080--1095. \url{https://doi.org/10.1080/00927872.2017.1339058}.
\bibitem{EreyDiamond} N. Erey, \emph{Powers of edge ideals with linear resolutions}, Comm. Algebra \textbf{46} (2018), no.~9, 4007--4020. \url{https://doi.org/10.1080/00927872.2018.1430810}.
\bibitem{EreyC4} N. Erey, \emph{Powers of ideals associated to $(C_4,2K_2)$-free graphs}, J. Pure Appl. Algebra \textbf{223} (2019), no.~7, 3071--3080. \url{https://doi.org/10.1016/j.jpaa.2018.10.009}.
\bibitem{EFFHKM} N. Erey, S. Faridi, T. H. H\`a, T. Hibi, S. Kara, and S. Morey, \emph{Gapfree graphs and powers of edge ideals with linear quotients}, Adv. in Appl. Math. \textbf{174} (2026), 103004. \url{https://doi.org/10.1016/j.aam.2025.103004}.
\bibitem{EHHKM} N. Erey, T. H. H\`a, T. Hibi, S. Kara, and S. Morey, \emph{Edge ideals of gapfree graphs with matching number $3$}, preprint, arXiv:2610.21866v1 (October 15, 2026). \url{https://arxiv.org/abs/2610.21866}.
\bibitem{Froberg} R. Fr\"oberg, \emph{On Stanley--Reisner rings}, in: Topics in Algebra, Part~2 (Warsaw, 1988), Banach Center Publ. \textbf{26}, PWN, Warsaw, 1990, 57--70.
\bibitem{HHZ} J. Herzog, T. Hibi, and X. Zheng, \emph{Monomial ideals whose powers have a linear resolution}, Math. Scand. \textbf{95} (2004), 23--32. \url{https://doi.org/10.7146/math.scand.a-14446}.
\bibitem{MV} N. C. Minh and T. Vu, \emph{A characterization of graphs whose small powers of their edge ideals have a linear free resolution}, Combinatorica \textbf{44} (2024), no.~2, 337--353. \url{https://doi.org/10.1007/s00493-023-00074-z}.
\bibitem{NP} E. Nevo and I. Peeva, \emph{$C_4$-free edge ideals}, J. Algebraic Combin. \textbf{37} (2013), 243--248. \url{https://doi.org/10.1007/s10801-012-0364-2}.
\end{thebibliography}
\end{document}
